Una partícula carregada és a sobre d'una taula horitzontal sense fricció i dins d'un camp elèctric homogeni i constant. La partícula està lligada amb un fil a un punt $P$ respecte del qual pot pivotar lliurement. Inicialment, la partícula està subjectada al punt $A$ i en repòs, de manera que el fil, que està tens, forma un angle de 60° respecte al camp elèctric.

\figura[0.4]{fig1.png}

\begin{apartat}{1,25}
En la figura anterior, representeu sobre la partícula la força elèctrica $\vec{F_e}$ i la força que fa el fil $\vec{T}$. Calculeu el mòdul de la força elèctrica que actua sobre la partícula quan és a la posició $A$. Aquesta força elèctrica serà constant al llarg de la trajectòria des de $A$ fins a $B$? Justifiqueu la resposta.
\end{apartat}

\begin{apartat}{1,25}
Calculeu el mòdul de la velocitat de la partícula quan passa pel punt $B$. Justifiqueu la resposta i indiqueu el principi físic en què us heu basat.
\end{apartat}

\dades{%
  $L = 1,0\ \mathrm{m}$.\\
  $m = 2,5\ \mathrm{g}$.\\
  $\vec{E} = 1,2\times 10^3\ \vec{\imath}\ \mathrm{V/m}$.\\
  $q = 3\ \mathrm{\mu C}$.}
